\section{Normed Products and Complex Pfaffians}
\label{sec:norm-products-pfaffians}

We apply \cref{thm:spectral-gradient-subcalibration} to octonionic products
and complex Pfaffians in every integer degree.
The multiplication maps over normed division algebras are harmonic morphisms; see
\cite[Section~4.5.2]{RiedlerThesis25}. The minimality of the real-part
Pfaffian cones also follows from the holomorphic construction of
Hoppe--Tkachev \cite[Corollary~5.2 and Remark~5.5]{HoppeTkachev19}.
Their smallest algebras have the spectral table, with Laplacian
coefficients \(32\) and \(24\), respectively. Both exceed
\(24-16/d\), giving the potential \(F|\nabla F|\).

\subsection{Products in a normed division algebra}

Let \(\mathbb K=\mathbb H\) or \(\mathbb O\), put
\(\delta=\dim_{\R}\mathbb K\), and fix an integer \(d\ge2\).
Choose a full parenthesization of the ordered product of \(d\)
variables, and write
\begin{equation}\label{eq:norm-product-definition}
 P:\mathbb K^d\longrightarrow\mathbb K,
 \qquad P(x_1,\ldots,x_d)=x_1\cdots x_d,
 \qquad F=\operatorname{Re}P.
\end{equation}
The fixed parentheses allow for the nonassociativity of
\(\mathbb O\). Set
\begin{equation}\label{eq:norm-product-generators}
 r_a=|x_a|^2,\qquad E_k=e_k(r_1,\ldots,r_d),\qquad
 R=E_1=\sum_a|x_a|^2,
\end{equation}
where \(e_k\) denotes the \(k\)-th elementary symmetric polynomial.
We use \(E_0=1\) and \(E_k=0\) for \(k>d\).

\begin{proposition}[Differential closure of normed products]
\label{prop:norm-product-closure}
For every parenthesization in \eqref{eq:norm-product-definition},
the polynomial \(F\) is real irreducible, homogeneous of degree \(d\), and
satisfies
\begin{equation}\label{eq:norm-product-F-closure}
\begin{gathered}
 \Delta F=0,\qquad \Gamma(F,F)=E_{d-1},\\
 \Gamma(F,E_k)=2(d-k+1)E_{k-1}F,
\end{gathered}
\end{equation}
\begin{equation}\label{eq:norm-product-E-closure}
\begin{aligned}
 \Delta E_k&=2\delta(d-k+1)E_{k-1},\\
 \Gamma(E_i,E_j)&=4\sum_{a=1}^d r_a
       \frac{\partial e_i}{\partial r_a}
       \frac{\partial e_j}{\partial r_a}.
\end{aligned}
\end{equation}
Consequently
\begin{equation}\label{eq:norm-product-minimality}
 \cL(F)=-2E_{d-2}F.
\end{equation}
The polynomial critical set is
\begin{equation}\label{eq:norm-product-critical-set}
 Z_d=\{\nabla F=0\}
 =\bigcup_{1\le a<b\le d}\{x_a=x_b=0\},
\end{equation}
and has real codimension \(2\delta\).
\end{proposition}

\begin{proof}
Multiplication in \(\mathbb K\) is real bilinear and satisfies
\(|uv|=|u||v|\); see
\cite[Chapter~1, Appendix, pp.~21--22]{Shapiro2000}.
A parenthesized product is therefore real
multilinear, and induction on its multiplication tree gives
\[
 |P(x)|^2=\prod_{a=1}^d r_a.
\]
With all variables except \(x_a\) fixed, the derivative
\(A_a=D_{x_a}P:\mathbb K\to\mathbb K\) is a linear map satisfying
\[
 |A_av|^2=\left(\prod_{b\ne a}r_b\right)|v|^2.
\]
By polarization, \(A_a^*A_a=(\prod_{b\ne a}r_b)I\), the usual
characterization of a linear similarity
\cite[Chapter~1, Lemma~1.3(1), pp.~15--16]{Shapiro2000}.
The source and target have the same dimension, so also
\(A_aA_a^*=(\prod_{b\ne a}r_b)I\); this remains true when the
coefficient is zero.  Since \(F=\ip{P}{1}\), it follows that
\[
 |\nabla_{x_a}F|^2=|A_a^*1|^2=\prod_{b\ne a}r_b.
\]
Summing over \(a\) proves \(\Gamma(F,F)=E_{d-1}\).
Every second derivative within a single variable block vanishes,
so \(\Delta F=0\).  The choice \(x_1=\cdots=x_d=1\) shows that
\(F\) is nonzero.

For irreducibility, a factorization \(F=AB\) must be homogeneous
in every variable block, by the lowest- and highest-degree argument
in \cref{lem:square-free-critical-set}. Since each block degree of
\(F\) is one, \(A\) and \(B\) depend on disjoint groups of blocks.
If both are nonconstant, retain one variable from each group and set
all remaining variables equal to \(1\). Every parenthesization then
reduces \(F\) to \(\operatorname{Re}(uv)=\langle u,\bar v\rangle\),
a nondegenerate bilinear form of rank \(\delta>1\)
\cite[Chapter~1, Appendix, pp.~21--22]{Shapiro2000}.
The specialized factorization would express it as a product of two
linear forms in separate variables, of rank one, a contradiction.

Euler's identity in each variable block gives
\(\Gamma(F,r_a)=2F\).  Hence
\[
 \Gamma(F,E_k)=2F\sum_a\frac{\partial e_k}{\partial r_a}
 =2(d-k+1)E_{k-1}F.
\]
Moreover, \(\Delta r_a=2\delta\),
\(\Gamma(r_a,r_b)=4r_a\delta_{ab}\), and each \(e_k\) is affine
in each individual \(r_a\).  These facts prove
\eqref{eq:norm-product-E-closure} by the chain rule.
Harmonicity and the identity for \(\Gamma(F,E_{d-1})\) now give
\eqref{eq:norm-product-minimality}.

Finally, the nonnegative sum
\(E_{d-1}=\sum_a\prod_{b\ne a}r_b\) vanishes precisely when at
least two blocks vanish.  Each largest linear component of this
set has dimension \(\delta(d-2)\), proving the codimension claim.
\end{proof}

The vector-valued map itself satisfies
\begin{equation}\label{eq:norm-product-harmonic-morphism}
 \Delta P=0,\qquad DP\,DP^*=E_{d-1}I_{\delta},\qquad
 DP(\nabla E_{d-1})=4E_{d-2}P.
\end{equation}
Indeed, the first two identities were established blockwise in the
proof, while multilinearity gives \(DP(\nabla r_a)=2P\), whence
the third.  Thus the target of the pullback defining our hypersurface
is the hyperplane \(\{\operatorname{Re}y=0\}\subset\mathbb K\).
The umbilic condition in \eqref{eq:norm-product-harmonic-morphism}
places its regular pullback within Riedler's umbilic framework
\cite[Theorem~5.1.4]{RiedlerThesis25}.  The direct identity
\eqref{eq:norm-product-minimality} additionally covers regular zeros
lying over \(P=0\), for instance points with exactly one zero block.

\begin{proposition}[Exact algebra and quotient image]
\label{prop:norm-product-exact-algebra}
For \(d\ge3\),
\begin{equation}\label{eq:norm-product-smallest-algebra}
 \cA_R(F)=\R[E_1,\ldots,E_d,F],
\end{equation}
and its \(d+1\) displayed generators are algebraically independent.
Their joint image is exactly
\begin{equation}\label{eq:norm-product-quotient-image}
 \left\{(e_1(r),\ldots,e_d(r),f):
             r\in[0,\infty)^d,\quad f^2\le e_d(r)\right\}.
\end{equation}
The spherical quotient dimension is \(d\).
For \(d=2\), the smaller algebra \(\cA_R(F)=\R[R,F]\) is already
closed and its spherical quotient dimension is one.
\end{proposition}

\begin{proof}
Apply \cref{prop:spectral-smallest-algebra} with
\(r_a=|x_a|^2\) and \(b_k=2\delta(d-k+1)>0\) from
\cref{prop:norm-product-closure}.

The norm identity implies \(F^2\le E_d\).
Conversely, choose an imaginary unit \(u\in\mathbb K\), and put
\[
 x_1=\sqrt{r_1}(\cos\theta+u\sin\theta),\qquad
 x_a=\sqrt{r_a}\quad(a\ge2).
\]
All bracketings then give \(F=\sqrt{e_d(r)}\cos\theta\).
This realizes every value in \eqref{eq:norm-product-quotient-image},
including its boundary.
At distinct positive \(r_a\), the Jacobian of
\((e_1,\ldots,e_d)\) is a nonzero Vandermonde determinant.
Taking \(\sin\theta\ne0\) adds a phase direction changing \(F\)
while fixing all \(E_k\).  The joint map thus has rank \(d+1\)
on a nonempty open set, which proves algebraic independence.
Fixing \(E_1=1\) lowers the rank by one.  In the quadratic case,
\(R,F\) have independent differentials at generic points and
give the asserted one-dimensional spherical quotient.
\end{proof}

In particular, different octonionic parenthesizations have the same
abstract differential closure and the same quotient image.  These
statements do not identify the corresponding embedded hypersurfaces
up to orthogonal congruence.

\begin{theorem}[Octonionic product cones in every degree]
\label{thm:octonionic-product-minimizing}
For every \(d\ge2\) and every fixed parenthesization, let
\(F=\operatorname{Re}(x_1\cdots x_d)\) on \(\mathbb O^d\).
Then the cone \(C_d=\{F=0\}\subset\R^{8d}\) is the support of the
locally mass-minimizing boundary \(\partial\llbracket\{F<0\}\rrbracket\).
A global two-sided subcalibration off the set \(Z_d\) in
\eqref{eq:norm-product-critical-set} is
\begin{equation}\label{eq:octonionic-product-potential}
 G=F\sqrt{E_{d-1}}=F|\nabla F|,
 \qquad X=\frac{\nabla G}{|\nabla G|}.
\end{equation}
Here \(Z_d\) has real codimension sixteen; it contains nonzero
points whenever \(d\ge3\). The cone is strictly area minimizing
and strictly stable, with
\begin{equation}\label{eq:octonionic-product-Hardy-gap}
 Q_{C_d}(u)\ge\frac{(4d+1)^2}{4}\int_{C_d}r^{-2}u^2,
 \qquad u\in C_c^\infty((C_d)_{\reg}\setminus\{0\}).
\end{equation}
\end{theorem}

\begin{proof}
By \cref{prop:norm-product-closure}, \(\Delta E_{d-1}=32E_{d-2}\)
and \(Z_d\) has codimension sixteen. Since \(32>24-16/d\),
\cref{thm:spectral-gradient-subcalibration}, with \(N=8d\), gives
all assertions, including \eqref{eq:octonionic-product-Hardy-gap}.
\end{proof}

\begin{remark}[The quaternionic closure]
\label{rem:quaternionic-product-threshold}
Quaternionic products have the same algebra and minimality identities,
but \(\Delta S=16B\). The potential \(F\sqrt S\) works for \(d=2\)
and fails for \(d\ge3\): at equal positive roots its boundary
coefficient is \((-4+8/d)B<0\), and the phase construction above
realizes arbitrarily small nonzero \(F\) with those roots fixed.
\end{remark}

\subsection{Complex Pfaffians and their minor norms}

Let \(W_d\) be the space of complex skew-symmetric \(2d\)-by-\(2d\)
matrices, with the Euclidean metric
\begin{equation}\label{eq:pfaffian-metric}
 |Z|^2=\sum_{i<j}|z_{ij}|^2=\tfrac12\tr(ZZ^*),
 \qquad \dim_{\R}W_d=2d(2d-1).
\end{equation}
The independent coordinates are the complex entries \(z_{ij}\),
\(i<j\).  For \(d\ge2\), set
\begin{equation}\label{eq:pfaffian-generators}
 f(Z)=\operatorname{Pf}(Z),\qquad F=\operatorname{Re}f,
 \qquad E_k(Z)=\sum_{\substack{I\subset\{1,\ldots,2d\}\\|I|=2k}}
                 |\operatorname{Pf}(Z_I)|^2.
\end{equation}
Here \(Z_I\) is the principal submatrix on \(I\), and
\(\operatorname{Pf}(Z_\varnothing)=1\).
In particular, \(E_1=R=|Z|^2\) and \(E_d=|f|^2\).

\begin{proposition}[Pfaffian differential closure]
\label{prop:pfaffian-closure}
For \(F,E_k\) defined by \eqref{eq:pfaffian-generators}, \(F\) is
real irreducible and
\begin{equation}\label{eq:pfaffian-F-closure}
\begin{gathered}
 \Delta F=0,\qquad\Gamma(F,F)=E_{d-1},\\
 \Gamma(F,E_k)=2(d-k+1)E_{k-1}F,
\end{gathered}
\end{equation}
\begin{equation}\label{eq:pfaffian-Laplacian}
 \Delta E_k=4(d-k+1)(2d-2k+1)E_{k-1}.
\end{equation}
At every point there exist nonnegative numbers \(r_1,\ldots,r_d\)
such that \(E_k=e_k(r)\), and
\begin{equation}\label{eq:pfaffian-Gamma}
 \Gamma(E_i,E_j)=4\sum_{a=1}^d r_a
       \frac{\partial e_i}{\partial r_a}
       \frac{\partial e_j}{\partial r_a}.
\end{equation}
In particular, \(\cL(F)=-2E_{d-2}F\).
For \(d\ge3\),
\begin{equation}\label{eq:pfaffian-smallest-algebra}
 \cA_R(F)=\R[E_1,\ldots,E_d,F],
\end{equation}
with algebraically independent generators and spherical quotient
dimension \(d\).  For \(d=2\), \(\cA_R(F)=\R[R,F]\).
\end{proposition}

\begin{proof}
Complexification gives
\(2F=\operatorname{Pf}(Z)+\operatorname{Pf}(W)\) for independent
skew-symmetric matrices \(Z,W\). It is linear in \(z_{12}\), with
coefficient \(\operatorname{Pf}(Z_{\{3,\ldots,2d\}})\).
That coefficient is coprime to the constant term, since the latter
contains \(w_{12}w_{34}\cdots w_{2d-1,2d}\) with coefficient one.
Gauss' lemma proves irreducibility
\cite[Chapter~VIII, Section~5, pp.~395--398]{Knapp16}.

Write \(\omega=\sum_{i<j}z_{ij}e_i\wedge e_j\) in an orthonormal
basis of \(\mathbb C^{2d}\).  The coefficient of \(e_I\) in
\(\omega^k/k!\) is \(\operatorname{Pf}(Z_I)\).  Consequently
\begin{equation}\label{eq:pfaffian-exterior-norm}
 E_k=\left|\frac{\omega^k}{k!}\right|^2,
\end{equation}
which proves that all \(E_k\) are invariant under unitary congruence
\(Z\mapsto UZU^T\).
The skew-symmetric unitary normal form is
\[
 Z=U\operatorname{diag}(\sigma_1J,\ldots,\sigma_dJ)U^T,
 \qquad \sigma_a\ge0,\qquad
 J=\begin{pmatrix}0&1\\-1&0\end{pmatrix}.
\]
For a direct proof, let \(T(v)=Z\bar v\).  Skew-symmetry gives
\(T^2=-ZZ^*\) and
\(\ip{v}{Tw}=-\ip{w}{Tv}\), where the Hermitian inner product is
conjugate linear in its first argument.  Choose a unit eigenvector
\(v\) of \(ZZ^*\) with positive eigenvalue \(\sigma^2\), and put
\(w=-Tv/\sigma\).  These identities imply that \(v,w\) are
orthonormal, \(Tv=-\sigma w\), and \(Tw=\sigma v\).
Their orthogonal complement is \(T\)-invariant.  Induction gives
the displayed normal form, with zero blocks at the last step.
At this normal form only
principal minors which are unions of complete two-by-two blocks
can have nonzero Pfaffian.  Thus \(E_k=e_k(\sigma_1^2,\ldots,
\sigma_d^2)\).

To compute first derivatives, more generally use the block diagonal
matrix with complex entries \(z_aJ\).  A nonzero principal Pfaffian
is a product of its block entries; its derivative in any off-block
coordinate vanishes.  A principal Pfaffian which vanishes contributes
zero to the first derivative of its squared modulus.  Therefore the
gradient of \(E_k\) has only block components at this matrix, and
\[
 \frac{\partial E_k}{\partial z_a}
   =\bar z_a\,e_{k-1}(r_1,\ldots,\widehat r_a,\ldots,r_d),
 \qquad r_a=|z_a|^2.
\]
The real Euclidean inner product of these gradients gives
\eqref{eq:pfaffian-Gamma}.  Unitary congruence preserves the metric
\eqref{eq:pfaffian-metric}, so this calculation holds everywhere.

The polynomial \(f\) is holomorphic, hence \(\Delta F=0\).
Differentiating its Pfaffian expansion gives
\[
 \frac{\partial f}{\partial z_{ij}}
   =\pm\operatorname{Pf}(Z_{\{1,\ldots,2d\}\setminus\{i,j\}}).
\]
For a holomorphic polynomial on complex Euclidean space,
\[
 |\nabla\operatorname{Re}f|^2=\sum_{i<j}|\partial f/\partial z_{ij}|^2.
\]
This proves \(\Gamma(F,F)=E_{d-1}\).
At the block diagonal matrix, complex bilinear extension of
\(\Gamma\) gives
\[
 \Gamma(f,E_k)
 =2\sum_a\frac{\partial f}{\partial z_a}
                \frac{\partial E_k}{\partial\bar z_a}
 =2(d-k+1)E_{k-1}f.
\]
Unitary congruence preserves every \(E_k\), and \(f(UZU^T)=\det(U)f(Z)\).
Hence the identity is covariant under this action.
Taking real parts proves the remaining identity in
\eqref{eq:pfaffian-F-closure}; polynomial continuity includes all
degenerate normal forms.

For the Laplacian, apply
\(\Delta|p|^2=4\sum_{i<j}|\partial p/\partial z_{ij}|^2\)
to each holomorphic principal Pfaffian in \eqref{eq:pfaffian-generators}.
Each derivative deletes a pair of indices.  A fixed set of \(2k-2\)
indices is obtained in exactly
\(\binom{2d-2k+2}{2}=(d-k+1)(2d-2k+1)\) ways, proving
\eqref{eq:pfaffian-Laplacian}.
The minimality identity follows from
\(-\tfrac12\Gamma(F,E_{d-1})=-2E_{d-2}F\).

Now apply \cref{prop:spectral-smallest-algebra} with
\(r_a=\sigma_a^2\) and \(b_k=4(d-k+1)(2d-2k+1)>0\).
To verify independence, take distinct positive \(r_a\), with
\(z_1=\sqrt{r_1}e^{i\theta}\) and \(z_a=\sqrt{r_a}\) for \(a\ge2\).
Then \(E_k=e_k(r)\) and \(F=\sqrt{e_d(r)}\cos\theta\).
The Vandermonde and phase argument in
\cref{prop:norm-product-exact-algebra} applies without change.
It also proves that the joint image is exactly
\eqref{eq:norm-product-quotient-image}.
\end{proof}

\begin{theorem}[Area minimization for complex Pfaffians]
\label{thm:pfaffian-minimizing}
For every \(d\ge2\), the cone
\[
 C_d=\{\operatorname{Re}\operatorname{Pf}(Z)=0\}
 \subset W_d\cong\R^{2d(2d-1)}
\]
is the support of the locally mass-minimizing boundary of
\(\{\operatorname{Re}\operatorname{Pf}(Z)<0\}\).
A two-sided subcalibration is the normalized gradient of
\begin{equation}\label{eq:pfaffian-potential}
 G=F\sqrt{E_{d-1}}=F|\nabla F|,
\end{equation}
defined off
\begin{equation}\label{eq:pfaffian-critical-set}
 Z_d=\{Z:\operatorname{rank}_{\mathbb C}Z\le2d-4\}
     =\{\nabla F=0\}.
\end{equation}
This exceptional set has real codimension twelve.
The cone is strictly area minimizing and strictly stable, with
\begin{equation}\label{eq:pfaffian-Hardy-gap}
 Q_{C_d}(u)\ge\frac{(4d^2-6d+1)^2}{4}\int_{C_d}r^{-2}u^2,
 \qquad u\in C_c^\infty((C_d)_{\reg}\setminus\{0\}).
\end{equation}
\end{theorem}

\begin{proof}
By \cref{prop:pfaffian-closure}, the spectral table holds with \(S=E_{d-1},\; B=E_{d-2},\; \Delta S=24B\),
and \(24\ge24-16/d\).

The normal form shows that \(S=0\) precisely when at most
\(d-2\) of the \(\sigma_a\) are nonzero, proving
\eqref{eq:pfaffian-critical-set}. The stratum of
complex skew-symmetric \(n\)-by-\(n\) matrices of rank \(2q\)
has complex dimension \(2q(n-2q)+\binom{2q}{2}\).
Indeed, choose its kernel, of dimension \(n-2q\), and then choose
a nondegenerate alternating form on the \(2q\)-dimensional quotient.
Its complex codimension is therefore
\(\binom{n-2q}{2}\).
For \(n=2d\) and \(q=d-2\), this is six; lower rank strata have
larger codimension.  Hence \(Z_d\) is a closed algebraic set of
real codimension twelve and has locally zero
\(\cH^{2d(2d-1)-1}\)-measure.

\Cref{thm:spectral-gradient-subcalibration}, with \(N=2d(2d-1)\),
gives the minimizing assertions and \eqref{eq:pfaffian-Hardy-gap}.
Here \(4d^2-6d+1>0\) for \(d\ge2\).
\end{proof}

